\scanpage{21}
\mspage{20}

Ainsi on trouve donc une application
\[
 \unsure{\mathcal{F}|\mathscr{F}}\longrightarrow\mathrm{Fig}(M),\quad F\longmapsto\unsure{\mathrm{Multomb}|\mathrm{Multiomb}}(F)
 \tag*{$(*)$}
\]
Elle est injective, car \inserted{$X$ est connu quand $\mathrm{Omb}(X)$ est connu (c'est le plus grand élément de $\mathrm{Omb}(X)$ pour $\ll$), donc $\widetilde F$ est connu quand on \unsure{connaît|connait} $\unsure{\mathrm{Multom}|\mathrm{Multiomb}}(F)$, donc $F$ aussi.}

Il est immédiat \unsure{que}
\[
 \deleted{F}G\leqslant\deleted{G}F\quad
 \unsure{\Longleftrightarrow|\Longrightarrow}\quad
 \underbrace{
 \underbrace{\mathrm{Multomb}(\deleted{F}G)}_{
 \substack{\text{les strates sont}\\\text{les }\mathrm{Omb}(\deleted{X}Y)\\\deleted{X}Y\in\widetilde{\deleted{F}G}}}
 \leqslant
 \underbrace{\mathrm{Multomb}(\deleted{G}F)}_{
 \substack{\text{les strates}\\\text{sont les}\\\mathrm{Omb}(X)\\X\in\widetilde F}}}_{
 \substack{\text{or }\widetilde{\unsure{G|C}}\subset\widetilde F\text{ est}\\
 \text{une partie}\\\text{fermée de }\widetilde F\\\text{si }G\leq F}}
 \tag*{(1.38)}
\]

\uline{19 \unsure{Juin|juin}}\qquad Finalement, j'ai changé de point de vue pour le sens géométrique de la relation $\deleted{F}G\ll\deleted{G}F$, que je \unsure{concevais} comme ``$G$ est une \unsure{sous figure|sousfigure} d'une subdivision de $F$''. Déjà dans \deleted{\unsure{les}} le cas simplicial (i.e la structure \unsure{combinatoire} des \deleted{\unsure{tout|tous}} multistrates étant \deleted{celle} une structure de simplexe : \unsure{l'ens.|ens.} ordonné des sous-ens d'un ens. fini \ldots), cela n'a pas l'air naturel. Par exemple, si on a une théorie modérée, et deux simplexes modérés $X,Y$ dans un $\mathbb{R}^{\unsure{n}}$ disons, on dira que $\deleted{X}Y\ll X$ si c'est le cas pour les figures ensemblistes correspondantes -- ce qui n'implique sans doute pas l'existence d'une \unsure{subd.|subdiv.} de $X$ telle que $Y$ en soit une multistrate \unsure{et|de} la
